This appendix collects the convention-sensitive choices used in
Sections~\ref{sec:sequential-histories}--\ref{sec:affine-cocycles}.  The choices
are not mathematical assertions, but every composition, sign, and coordinate
formula depends on them.

\subsection{Chronological words and functional composition}

\begin{convention}[Order inside a history]
\label{conv:history-order}
A history
\[
  \gamma=(g_1,\ldots,g_n)
\]
performs $g_1$ first and $g_n$ last.  Its functional evaluation is therefore
\[
  \ev(\gamma)=g_n\circ\cdots\circ g_1
\]
when the $g_i$ are viewed as maps.
\end{convention}

\begin{convention}[Chronological concatenation]
\label{conv:chronological-concatenation}
The written product $\gamma\delta$ means first perform $\gamma$ and then
perform $\delta$.  Thus
\[
  \ev(\gamma\delta)=\ev(\delta)\circ\ev(\gamma).
\]
The notation is chronological rather than functional: the earlier process is
written on the left.
\end{convention}

For example, if $g(z)=z+\mu$ and $h(z)=qz$, then the word
$\gamma=(g,h)$ and the chronological product of the one-step histories $gh$
both act as
\[
  h\circ g(z)=qz+q\mu.
\]
The reverse chronological word $hg$ acts as $g\circ h(z)=qz+\mu$.  This
two-step check fixes the sign of their endpoint difference:
\[
  \ev(gh)(z)-\ev(hg)(z)=\mu(q-1).
\]

\subsection{Column-vector projective action}

\begin{convention}[Matrix action]
\label{conv:column-projective-action}
Homogeneous coordinates are columns, and matrices act from the left:
\[
  \begin{pmatrix}A&B\\C&D\end{pmatrix}
  \begin{pmatrix}u\\v\end{pmatrix}
  =\begin{pmatrix}Au+Bv\\Cu+Dv\end{pmatrix}.
\]
On the affine chart $z=[z:1]$, this gives
\[
  M\cdot z=\frac{Az+B}{Cz+D}.
\]
\end{convention}

If $M_i$ represents $g_i$, the history $(g_1,\ldots,g_n)$ is represented by
\[
  M_n\cdots M_1.
\]
Consequently, for chronological concatenation,
\begin{equation}
  \rho(\gamma\delta)=\rho(\delta)\rho(\gamma).
  \label{eq:appendix-rho-composition-order}
\end{equation}
This is the matrix version of
Convention~\ref{conv:chronological-concatenation}.  A row-vector right action
would require a different product convention; it is not used in this paper.

Projective matrices are scalar classes.  If $M$ and $aM$ differ by
$a\in K^\times$, they represent the same element of
$\operatorname{PGL}_2(K)$.  Matrix identities used for projective maps may
therefore be checked either exactly in $\operatorname{GL}_2(K)$ or up to one
overall nonzero scalar, provided the choice is stated.

\subsection{Ordinary and projective domains}

The following table records the domain distinction for the contexts involving
division.  Here $c\in K$ and $z$ is the evolving affine value.

\begin{center}
\renewcommand{\arraystretch}{1.35}
\begin{tabular}{llll}
\toprule
Context & Ordinary rule & Ordinary condition & Projective status \\
\midrule
$C_{/,c}^{(1)}[z]=z/c$
& divide by fixed $c$
& $c\ne0$
& invertible if $c\ne0$ \\
$C_{/,c}^{(2)}[z]=c/z$
& divide by evolving $z$
& $z\ne0$
& invertible if $c\ne0$ \\
$C_{\times,c}^{(\varepsilon)}[z]=cz$
& multiply by $c$
& none
& invertible if $c\ne0$ \\
\bottomrule
\end{tabular}
\end{center}

Two different exclusions appear in the table.  The parameter condition
$c\ne0$ ensures that the context has an invertible projective matrix.  The
state condition $z\ne0$ ensures that an ordinary second-slot division is
defined at that step.  Projectively, a non-degenerate $c/z$ sends $0$ to
$\infty$; this is a chart transition, not an ordinary quotient in $K$.

For a composite history, ordinary admissibility is checked at every
intermediate state.  It is not enough that the final projective point lies in
the affine chart.  In particular, two projective inversions compose to the
identity operator, but the corresponding ordinary two-step word is not
admissible at zero.

\subsection{Left and right group data}

The affine group is represented by
\[
  g(\Phi,\xi)=
  \begin{pmatrix}\Phi&\xi\\0&1\end{pmatrix},
  \qquad \Phi\in K^\times.
\]
With the column action, matrix multiplication gives
\begin{equation}
  g(\Phi_2,\xi_2)g(\Phi_1,\xi_1)
  =g(\Phi_2\Phi_1,\Phi_2\xi_1+\xi_2).
  \label{eq:appendix-affine-group-law}
\end{equation}
Thus the earlier translation is rescaled by the later affine map.  Equation
\eqref{eq:appendix-affine-group-law} is the two-factor origin of the
target-frame cocycle.

Over the positive real component, write $\Phi=e^M$.  The two
Maurer--Cartan forms used in Section~\ref{sec:affine-cocycles} have the
following readings.

\begin{center}
\begin{tabular}{lll}
\toprule
Form & Translation component & Frame \\
\midrule
$g^{-1}dg$ & $e^{-M}d\xi$ & body/source \\
$dg\,g^{-1}$ & $d\xi-\xi dM$ & spatial/target \\
\bottomrule
\end{tabular}
\end{center}

The words ``left'' and ``right'' here refer to invariant differential forms on
the affine group.  They do not refer to first and second operand slots.  The
operand-slot mirror is syntactic; the Maurer--Cartan distinction is
group-theoretic.

\subsection{Target and source-normalized translation coordinates}

For $x\mapsto\Phi x+\xi$, the symbol $\xi$ always denotes target-frame
translation.  The source-normalized coordinate is
\[
  \widehat\xi=\Phi^{-1}\xi.
\]
Their discrete and differential conversions are
\[
  \xi=\Phi\widehat\xi,
  \qquad
  \widehat\xi=\Phi^{-1}\xi,
\]
and, when $\Phi=e^M$,
\[
  d\xi=\xi\,dM+e^M d\widehat\xi.
\]
No formula in the main text changes from $\xi$ to $\widehat\xi$ without this
normalization being displayed.

\subsection{Mirror, reversal, and inverse notation}

For $\gamma=(g_1,\ldots,g_n)$, the three operations are summarized by
\begin{center}
\begin{tabular}{lll}
\toprule
Operation & Resulting word & Required hypothesis \\
\midrule
Mirror $m\gamma$
& $(m(g_1),\ldots,m(g_n))$
& none \\
Temporal reversal $r\gamma$
& $(g_n,\ldots,g_1)$
& none \\
Path inverse $\gamma^{-1}$
& $(g_n^{-1},\ldots,g_1^{-1})$
& every $g_i$ invertible \\
\bottomrule
\end{tabular}
\end{center}
Mirror exchanges operand slots and retains time.  Temporal reversal changes
time and retains the original contexts.  Path inverse reverses time and
inverts every operator.  Accordingly,
\[
  \rho(\gamma^{-1})=\rho(\gamma)^{-1},
\]
whereas no analogous identity holds for $r\gamma$ in general.

\subsection{Notation conversion summary}

The table below lists older or potentially ambiguous descriptions and the
notation used in Paper~I.

\begin{center}
\renewcommand{\arraystretch}{1.25}
\begin{tabular}{ll}
\toprule
Potentially ambiguous description & Paper~I notation \\
\midrule
evolving value in operand slot one
& $C_{\omega,c}^{(1)}[z]=\omega(z,c)$ \\
evolving value in operand slot two
& $C_{\omega,c}^{(2)}[z]=\omega(c,z)$ \\
chronological context word
& $\gamma=(g_1,\ldots,g_n)$, $g_1$ first \\
history product
& $\gamma\delta$, first $\gamma$, then $\delta$ \\
operator evaluation
& $\rho(\gamma)$ \\
ordinary endpoint at $x$
& $\nu_x(\gamma)$ \\
accumulated scale
& $\Phi_\gamma$ \\
target-frame translation
& $\xi_\gamma$ \\
source-normalized translation
& $\widehat\xi_\gamma=\Phi_\gamma^{-1}\xi_\gamma$ \\
operand-slot exchange
& $m\gamma$ \\
time-order reversal
& $r\gamma$ \\
inverse operator path
& $\gamma^{-1}$ \\
\bottomrule
\end{tabular}
\end{center}
